\documentclass[conference]{IEEEtran}
\IEEEoverridecommandlockouts
% The preceding line is only needed to identify funding in the first footnote. If that is unneeded, please comment it out.
\usepackage{cite}
\usepackage{amsmath,amssymb,amsfonts}
\usepackage{algorithmic}
\usepackage{graphicx}
\usepackage{textcomp}
\usepackage{xcolor}
\usepackage{svg}
\usepackage{url}
\usepackage[hidelinks]{hyperref}
\def\BibTeX{{\rm B\kern-.05em{\sc i\kern-.025em b}\kern-.08em
    T\kern-.1667em\lower.7ex\hbox{E}\kern-.125emX}}
\begin{document}

\title{Summer Research Internship (SRI) at DAU on\\
{Generative Refinement for Target Speech Extraction using Joint Mixture-Guided Diffusion Priors}\\[1ex]
\large \textbf{Faculty Mentor: Dr. Hemant Patil}%
\thanks{This work was made possible by the SRI program at DAU.}
}

\author{\IEEEauthorblockN{Aayush Prajapati}
\IEEEauthorblockA{\textit{Electrical and Computer Science Engineering } \\
\textit{IITRAM}\\
Ahemdabad, India \\
Aayush.Prajapati.24co@iitram.ac.in\\
+91 9879401314}
}

\date{July 2026}

\maketitle

\begin{abstract}
Target Speech Extraction (TSE) conditioned on a speaker's reference cue isolates a desired speaker from a multi-talker mixture. While discriminative models excel at voice separation, they frequently introduce artificial processing distortions, yielding a robotic and unnatural output. Generative diffusion priors can restore natural speech characteristics, but suffer from speaker identity drift and speech hallucinations in overlapping or target-silent regions due to the lack of acoustic boundaries. In this work, we propose Joint Mixture Guidance, a training-free generative posterior sampling framework designed to resolve these limitations. By defining a dynamic energy-based target mask from the discriminative output, our approach guides the prior using a masked $L_2$ loss to match the mixture signal only in active target zones, preventing speech hallucinations in silent or overlapping regions. Evaluated on the Libri2Mix benchmark, our hybrid pipeline achieves a significant perceptual quality boost, raising the DNSMOS score by $+0.340$ under $T=400$ sampling steps and $+0.345$ under a low-latency $T=100$ configuration, while conserving target speaker similarity and word intelligibility.
\end{abstract}


\begin{IEEEkeywords}
Target Speech Extraction, Diffusion Posterior Sampling,  Joint Mixture Guidance
\end{IEEEkeywords}

\section{Introduction}
This report summarizes the design, implementation, and evaluation of a high-fidelity Target Speech Extraction system developed during the Summer Research Internship.

\subsection{Background}
The human auditory system possesses the remarkable ability to focus on a single target talker in a crowded environment, a phenomenon known as the cocktail party effect\cite{cherry}. In speech signal processing, Target Speech Extraction (TSE) aims to mimic this ability by isolating a target speaker's voice from a mixture, conditioned on a reference clue (such as a pre-recorded enrollment utterance). Unlike Blind Source Separation (BSS) which separates all active sources and suffers from permutation ambiguity, TSE extracts only the target source of interest, making it highly practical for applications in hearing aids, smart assistants, and automatic speech recognition (ASR) front-ends.

\subsection{Problem Statement: An Overview}Discriminative and generative models are two different approaches for modeling any signal processing tasks. Discriminative models are very good at isolating the target speaker and cutting out background noise. Because they are trained directly to separate the voices, they can clean up the audio and make the target speaker's words much easier to hear.

However, these models have a major downside: they often introduce digital artifacts, musical noise, and make the extracted voice sound robotic and unnatural. 

Generative models, such as diffusion-based priors (like ArrayDPS), solve this problem. Instead of just filtering the noisy audio, they try to reconstruct the natural human voice from scratch. By starting with the output of the discriminative model, the diffusion prior can restore the natural pitch, tone, and high-frequency details of the speaker, making the audio sound much more organic.

But standard generative diffusion models have their own failure modes. Because they run unconditionally, they do not know what the original mixture sounded like. This leads to two major issues:
\begin{itemize}
    \item \textbf{Speaker Identity Drift:} The model can lose track of the target speaker's voice and start morphing it into a different person's voice.
    \item \textbf{Speech Hallucinations:} In silent gaps or overlapping speech, the model starts ``hallucinating'' and generating random words that were never spoken, which ruins the accuracy of the transcription (raising the Word Error Rate).
\end{itemize}
To fix these issues, we propose \textbf{Joint Mixture Guidance}. It sets a physical limit on the diffusion process using the original mixture, preventing the model from hallucinating or changing the speaker's voice while still keeping the audio natural.





\subsection{Related Works}
Our project builds on several key papers in the target speech extraction and generative audio space:
\begin{itemize}
    \item \textbf{NeuralTSE Survey \cite{b1}:} We adopted this comprehensive review to understand the fundamentals of target speech extraction. It helped us map out how different speaker clues (like enrollment audio) are encoded and fused with mixture signals.
    \item \textbf{ArrayDPS \cite{b2}:} This paper showed how to use a pre-trained single-speaker diffusion model as a prior to separate multi-speaker mixtures without needing array geometry info. We use a similar 1D U-Net prior as our generative backbone.
    \item \textbf{SoloSpeech \cite{b4}:} This is a recent generative TSE model that uses latent diffusion to extract voices. While it works well, it is highly complex and requires training a custom time-frequency VAE from scratch. However, we adopted their comprehensive evaluation framework (incorporating SI-SNR, PESQ, ESTOI, DNSMOS, WER, and speaker similarity) to benchmark our model across multiple perceptual dimensions.
    \item \textbf{Uni-ArrayDPS \cite{b3}:} Similar to our approach, this work uses a diffusion prior to refine the output of a discriminative model without retraining. However, it focuses on using spatial covariance matrices to whiten noise in multi-channel microphone arrays, whereas our work focuses on bounding the signal envelope in single-channel mixtures.
    \item \textbf{TF-GridNet \cite{b5}:} A high-performance time-frequency domain discriminative model that serves as our baseline separation front-end.
\end{itemize}

\subsection{Contributions}
The contributions of this internship project are as follows:
\begin{enumerate}
    \item We propose \textbf{Joint Mixture Guidance (V1)}, a training-free posterior sampling framework that prevents speech hallucinations in generative priors.
    \item We employ a speaker-independent \textbf{USEF-TSE\cite{b6} (TFGridNet)} front-end and resolve its computational and memory limits on long audio mixtures by implementing a sliding-window audio chunking mechanism.
    \item We design an energy-based \textbf{target mask} from the discriminative output to separate clean target speech from silent or overlapping sections.
\end{enumerate}

\subsection{Notations}
Let $y_{\text{mix}}(t)$ denote the time-domain mixture signal, and $s(t)$ represent the clean target speech. The estimated output from the discriminative network is represented as $y_{\text{disc}}(t)$. The target speaker's energy mask is denoted as $M$. The reverse diffusion state at step $t$ is $x_t$, and its denoised estimate is $x_0$.

\subsection{Overview of the Report}
The rest of this report is organized as follows. Section II describes the System Model and formulates the discriminative target speech extraction network. Section III details the Proposed Approach, presenting the mathematical derivation of the Joint Mixture-Guided generative refinement framework. Section IV presents the Numerical Results, including the experimental setup, evaluation metrics, and performance comparisons. Section V concludes the report and outlines Future Work.

\section{System Model}
This section defines the mathematical layout of our speech mixture model.

\subsection{Signal Model}
\begin{equation}
y_{\text{mix}}(t) = s_1(t) + s_2(t) + n(t), \label{eq:mixture}
\end{equation}
where $s_1(t)$ and $s_2(t)$ are the clean speech signals of Speaker 1 and Speaker 2 respectively, and $n(t)$ represents additive noise. Given a reference enrollment cue $c_s$ for the target speaker $s \in \{1, 2\}$, the TSE network $f_{\phi}$ extracts an initial approximation:
\begin{equation}
y_{\text{disc}}(t) = f_{\phi}(y_{\text{mix}}(t), c_s). \label{eq:disc}
\end{equation}

\subsection{Detailed Problem Statement}
While $y_{\text{disc}}$ preserves the lexical structure of the target speaker, it contains high-frequency processing artifacts. The goal of the diffusion prior is to sample from the posterior distribution $p(x | y_{\text{mix}}, y_{\text{disc}})$. During overlapping intervals, the energy mask $M(t, f)$ computed from $y_{\text{disc}}$ contains leaking energy from the interfering speaker. This leakage leads to a trade-off: if the generative prior is allowed to inpaint these regions without bounds, it enhances the interference, leading to cross-speaker hallucinations and speaker identity drift.

\section{Proposed Approach}
We present our Joint Mixture Guidance framework below.

\subsection{A Mathematical Derivation}
Our generative refinement is based on a pre-trained single-speaker diffusion model called ArrayDPS\cite{b1}. In a diffusion model, the network learns to clean up noisy audio step-by-step. To initialize the process, we take the output of our discriminative model ($y_{\text{disc}}$), inject a small amount of noise, and run the reverse diffusion loop to restore natural speech details.

To make sure the diffusion model stays accurate to the target speaker and doesn't make up random sounds, we guide the denoising process at each step. We do this by calculating a guidance loss $\mathcal{L}_{\text{guidance}}$ and using its gradient to adjust the denoising trajectory:
\begin{equation}
\nabla_{x_t} \log p_t(x_t | y_{\text{mix}}, y_{\text{disc}}) \approx s_{\theta}(x_t, t) - \nabla_{x_t} \mathcal{L}_{\text{guidance}}(x_0),
\end{equation}

The overall architectural pipeline of this generative posterior sampling framework is illustrated in Fig.~\ref{fig:architecture}.
\begin{figure*}[htbp]
\centerline{\includesvg[width=0.9\textwidth]{Proposed_Approch_drawio.svg}}
\caption{Block diagram of the proposed joint mixture-guided target speech extraction framework.}
\label{fig:architecture}
\end{figure*}


where $s_{\theta}(x_t, t) is the standard denoiser score, and $x_0$ is the estimated clean speech at the current step.

Before calculating the guidance loss, we estimate where the target speaker is active. We compute the local signal energy $E_y(t)$ of the discriminative output $y_{\text{disc}}(t)$ and apply a threshold to get a binary target mask $M(t)$:
\begin{equation}
M(t) = \mathbb{I}\left(E_y(t) \geq \gamma \cdot \max_t E_y(t)\right),
\end{equation}
where $\gamma = 0.005$ is the energy threshold parameter. The mask is $1$ in active target speech regions and $0$ in silent or overlapping regions.

We then define our guidance loss $\mathcal{L}_{\text{guidance}}$ as the masked $L_2$ norm between the predicted clean speech $x_0$ and the original mixture $y_{\text{mix}}$:
\begin{equation}
\mathcal{L}_{\text{guidance}} = \| M \odot (x_0 - y_{\text{mix}}) \|_2.
\end{equation}
By using the target speaker mask $M$, we force the diffusion model to align with the mixture energy only in active target regions, while allowing the prior to freely denoise and refine silent or overlapping intervals where $M = 0$.

\subsection{Salient Aspects of the Proposed Solution}
Our proposed solution has three salient properties:
\begin{itemize}
    \item \textbf{Training-Free Adaptability:} The framework does not require retraining or joint training of either network, meaning any high-performance discriminative model can be easily integrated.
    \item \textbf{Dynamic Target Gating:} The energy-based mask $M(t)$ automatically partitions the audio, deciding where to track the discriminative speech and where to bound the prior to prevent hallucinations.
    \item \textbf{Single-Channel Compatibility:} Unlike multi-channel refinement systems that require microphone arrays and spatial covariance matrices, our pipeline operates on a single channel by utilizing the mixture envelope directly as a physical boundary constraint.
\end{itemize}

\subsection{Assumptions behind the Proposed Solution}
The system operates under three fundamental assumptions:
\begin{itemize}
    \item The discriminative output $y_{\text{disc}}$ preserves the correct target words, acting as a reliable structural anchor.
    \item The target speaker's voice amplitude is physically bounded by the mixture envelope, meaning the separated speaker cannot be louder than the original mix at any point.
    \item The pre-trained diffusion model is trained on clean speech, allowing it to guide the noisy discriminative output towards clean speech statistics.
\end{itemize}

\subsection{Conceptual Insights}
Combining discriminative extraction and generative refinement bridges the gap between signal-to-noise ratio (SNR) and natural human timbre. Discriminative models (TFGridNet) are good at global voice extraction but create artificial phase distortions. Generative models (ArrayDPS) generate natural, organic phases but tend to drift or hallucinate words in overlaps. Gating the guidance prevents these hallucinations, keeping the generated target speaker physically inside the mixture.

\subsection{Analytical Comparison with Related Works}
We compare our Joint Mixture Guidance method against two alternative baselines:
\begin{itemize}
    \item \textbf{Standard DPS:} Standard DPS does not utilize a mixture boundary. Because it has no reference to the mixture in silent regions, it drifts and has no baseline to anchor, causing word hallucinations in pauses.
    \item \textbf{Masked DPS:} Masked DPS disables guidance in silent gaps. While this works in silent pauses, it fails in overlapping speech because the mask cannot separate the target speaker from the background speaker. This leads to the prior enhancing the background speaker instead.
\end{itemize}
Our method applies the mask to the original mixture and uses its envelope as a ceiling, preventing background speakers from leaking through during overlaps.

\section{Numerical Results}
We evaluate our V1 Joint Mixture Guidance pipeline on the Libri2Mix\footnote{\url{https://github.com/JorisCos/LibriMix}}~\cite{librimix} 8kHz test dataset. 

\subsection{Evaluation Metrics}
We assess three key dimensions of the extracted speech:
\begin{enumerate}
    \item \textbf{Intrusive Quality:} We use Scale-Invariant Signal-to-Noise Ratio (SI-SNR)~\cite{sisnr} and Extended Short-Time Objective Intelligibility (ESTOI)~\cite{estoi} to evaluate waveform alignment and intrusive speech quality.
    \begin{table*}[t]
\caption{Target speech extraction results on the Libri2Mix dataset. $\star$ are the results of models reproduced by us. D/G indicates if the model is discriminative or generative. Bold for the best result and \underline{underline} for the second-best result.}
\label{tab:results}
\begin{center}
\begin{tabular}{|l|c|cc|c|cc|}
\hline
\textbf{Method} & \textbf{D/G} & \multicolumn{2}{c|}{\textbf{Perceptual Quality}} & \textbf{Naturalness} & \multicolumn{2}{c|}{\textbf{Intelligibility}} \\
 & & \textbf{SI-SNR} $\uparrow$ & \textbf{ESTOI} $\uparrow$ & \textbf{DNSMOS} $\uparrow$ & \textbf{WER (\%)} $\downarrow$ & \textbf{SIM} $\uparrow$ \\
\hline
\textit{T = 100 Configurations:} & & & & & & \\
Raw Mixture & -- & -0.128 & 0.538 & 2.939 & N/A & N/A \\
TFGridNet $\star$ & D & \textbf{12.450} & \textbf{0.877} & \underline{3.201} & \textbf{1.979} & \textbf{0.982} \\
Proposed (Our V1) $\star$ & G & \underline{12.287} & \underline{0.834} & \textbf{3.546} & \underline{2.445} & \underline{0.978} \\
\hline
\textit{T = 400 Configurations:} & & & & & & \\
Raw Mixture & -- & -0.235 & 0.550 & 3.006 & N/A & N/A \\
TFGridNet $\star$ & D & \textbf{12.341} & \textbf{0.865} & \underline{3.209} & \underline{1.662} & \textbf{0.981} \\
Proposed (Our V1) $\star$ & G & \underline{12.111} & \underline{0.825} & \textbf{3.549} & \textbf{1.512} & \underline{0.978} \\
\hline
\textit{Noisy Environment (T = 100):} & & & & & & \\
Raw Noisy Mix & -- & -1.390 & 0.418 & 2.126 & N/A & N/A \\
TFGridNet $\star$ & D & \textbf{9.828} & \textbf{0.768} & \underline{3.104} & \textbf{14.466} & \underline{0.955} \\
Proposed (Our V1) $\star$ & G & \underline{9.762} & \underline{0.742} & \textbf{3.321} & \underline{16.369} & \textbf{0.965} \\
\hline
\end{tabular}
\end{center}
\end{table*}
    \item \textbf{Perceptual Quality and Naturalness:} We utilize Microsoft's DNSMOS~\cite{dnsmos} using offline PyTorch implementation from DNSMOSPro\footnote{\url{https://github.com/fcumlin/DNSMOSPro}}~\cite{dnsmospro} to estimate the overall perceptual quality and naturalness of the audio files.

    
    
    \item \textbf{Intelligibility and Speaker Identity:} For the intelligibility metrics, we performed ASR on the speech files using the Whisper large-v3-turbo model\footnote{\url{https://huggingface.co/openai/whisper-large-v3-turbo}}~\cite{whisper} and calculated the Word Error Rate (WER) using the \texttt{jiwer} toolkit\footnote{\url{https://github.com/jitsi/jiwer}}. In addition, we measured the cosine similarity between the estimated speech and the ground truth using a pre-trained WavLM-based speaker verification network\footnote{\url{https://huggingface.co/microsoft/wavlm-base-plus-sv}}~\cite{wavlm} to assess speaker similarity (SIM).
\end{enumerate}

\subsection{Performance Comparison}
Table~\ref{tab:results} shows the target speech extraction evaluation results across the different models on the Libri2Mix dataset.

As expected, generative refinement (Proposed) introduces a minor reduction in waveform-matching metrics (SI-SNR and ESTOI) compared to the discriminative baseline (TFGridNet), as generative priors synthesize natural speech details rather than strictly matching the noisy source waveform. However, this drop is negligible ($\leq 0.230$~dB in SI-SNR across all clean test conditions). 

Meanwhile, as per our primary design goal, the perceptual quality and naturalness (DNSMOS) are dramatically improved—boosting the score by $+0.345$ for $T=100$ steps and $+0.340$ for $T=400$ steps. Intelligibility is highly preserved; in fact, the WER improves as the diffusion process is run for more steps, yielding a $0.150\%$ absolute reduction under $T=400$ compared to the discriminative baseline. Furthermore, under noisy conditions, the generative prior improves target speaker similarity (SIM) by $+0.010$ over TFGridNet, acting as a stabilizer for the speaker's vocal fingerprint.


\section{Conclusion}

\subsection{Summary}
In this report, we presented a training-free hybrid Target Speech Extraction pipeline incorporating Joint Mixture Guidance. By guiding the generative diffusion prior using an energy-based target mask derived from the discriminative network (TFGridNet), our approach resolves the vocal hallucinations and speaker identity drift typical of unconditional diffusion sampling. The framework achieves a massive naturalness boost of up to $+0.345$ DNSMOS, while maintaining high speaker similarity ($0.978$) and word intelligibility. Furthermore, the generative prior acts as a speaker identity stabilizer in noisy environments, improving speaker similarity by $+0.010$.

\subsection{Future Work}
Future work will focus on:
\begin{itemize}
    \item Testing a ReLU-based amplitude boundary in the guidance loss to enforce physical mixture limits.
    \item Replacing the 1D U-Net in the diffusion step with speech-specific networks like WaveNet.
    \item Fine-tuning the discriminative front-end and the diffusion prior jointly instead of keeping them training-free.
\end{itemize}


\section*{Acknowledgment}
The author would like to thank the faculty mentor and the staff and members of the Speech Research Laboratory (SRL) at DAU for their invaluable support, computing resources, and guidance throughout this project. The author also thanks the Summer Research Internship (SRI) program for providing this opportunity.

\begin{thebibliography}{00}
\bibitem{cherry} E. C. Cherry, ``Some Experiments on the Recognition of Speech, with One and with Two Ears,'' \textit{The Journal of the Acoustical Society of America}, vol. 25, no. 5, pp. 975--979, 1953.
\bibitem{b2} Z. Xu, X. Fan, Z.-Q.王, X. Jiang, and R. R. Choudhury, ``ArrayDPS: Unsupervised Blind Speech Separation with a Diffusion Prior,'' in \textit{Proceedings of the International Conference on Machine Learning (ICML)}, 2025.
\bibitem{b1} K. Zmolikova, M. Delcroix, T. Ochiai, K. Kinoshita, J. Cernocky, and D. Yu, ``Neural Target Speech Extraction: An Overview,'' \textit{IEEE Journal}, 2023.
\bibitem{b4} H. Wang, J. Hai et al., ``SoloSpeech: Enhancing Intelligibility and Quality in Target Speech Extraction through a Cascaded Generative Pipeline,'' \textit{IEEE Journal}, 2025.
\bibitem{b3} Z. Xu, A. Pandey, J. Azcarreta et al., ``Uni-ArrayDPS: Unified Diffusion Refinement for Multi-Channel Speech Enhancement and Separation,'' \textit{IEEE Journal}, 2025.
\bibitem{b5} Z.-Q. Wang, S. Cornell, S. Choi, Y. Lee, B.-Y. Kim, and S. Watanabe, ``TF-GridNet: Making Time-Frequency Domain Models Great Again for Monaural Speaker Separation,'' \textit{IEEE/ACM Transactions on Audio, Speech, and Language Processing}, 2023.
\bibitem{b6} Z.-Q. Wang et al., ``USEF-TSE: Unified Speaker Extraction Framework with Time-Frequency Domain Models,'' \textit{IEEE Journal}, 2024.
\bibitem{librimix} J. Cosentino, M. Pariente, S. Cornell, A. Deleforge, and E. Vincent, ``LibriMix: An Open-Source Dataset for Generalizable Speech Separation,'' in \textit{Proceedings of the International Conference on Association for Computational Linguistics (Interspeech)}, 2020.
\bibitem{sisnr} J. Le Roux, S. Wisdom, H. Erdogan, and J. R. Hershey, ``SDR--half-baked or well done?'' in \textit{IEEE International Conference on Acoustics, Speech and Signal Processing (ICASSP)}, 2019, pp. 626--630.
\bibitem{estoi} C. H. Taal, R. C. Hendriks, R. Heusdens, and J. Jensen, ``An Algorithm for Intelligibility Prediction of Time--Frequency Weighted Noisy Speech,'' \textit{IEEE Transactions on Audio, Speech, and Language Processing}, vol. 19, no. 7, pp. 2125--2136, 2011.
\bibitem{dnsmos} C. K. Reddy, V. Gopal, and R. Cutler, ``DNSMOS: A Non-Intrusive Complexity-Effective Metric for Benchmark Evaluation of Speech Quality in Noise,'' in \textit{ICASSP}, 2021.
\bibitem{dnsmospro} F. Cumlin, ``DNSMOSPro: A PyTorch-based Local Implementation of Microsoft DNSMOS,'' GitHub Repository, 2023. [Online]. Available: https://github.com/fcumlin/DNSMOSPro
\bibitem{whisper} A. Radford, J. W. Kim, T. Xu, G. Brockman, C. McLeavey, and I. Sutskever, ``Robust Speech Recognition via Large-Scale Weak Supervision,'' in \textit{ICML}, 2023.
\bibitem{wavlm} S. Chen, C. Wang, Z. Chen, Y. Wu, S. Liu, J. Chen, J. Li, N. Kanda, T. Yoshioka, and X. Yu, ``WavLM: Large-Scale Self-Supervised Pre-Training for Full Stack Speech Processing,'' \textit{IEEE Journal of Selected Topics in Signal Processing}, vol. 16, no. 6, pp. 1505--1518, 2022.
\bibitem{b7} M. Young, \textit{The Technical Writer's Handbook}. Mill Valley, CA: University Science, 1989.
\end{thebibliography}

\end{document}
